% Part: first-order-logic
% Chapter: beyond
% Section: second-order-logic

\documentclass[../../../include/open-logic-section]{subfiles}

\begin{document}

\olfileid{fol}{byd}{sol}

\olsection{د دويمې درجې منطق}

د دويمې درجې منطق ژبه يوازې د افرادو پر !!a{domain} کميت ټاکنې ته
اجازه نۀ ورکوي، بلکې د هغه !!{domain} پر اړيکو هم کميت ټاکل پرې کېدے
شي۔ د لومړۍ درجې يوې ژبې~$\Lang{L}$ ته د هر~$k$ دپاره داسې
!!{variable}s~$R$ ورزياتېږي چې پر $k$-ځاييزو اړيکو ګرځي، او پر دغو
متغيرونو کميت ټاکنه هم روا کېږي۔ که $R$ د يوې $k$-ځاييزې اړيکې
!!a{variable} وي او $t_1$، \dots،~$t_k$ عادي (د لومړۍ درجې) ترمونه
وي، نو $\Atom{R}{t_1,\dots,t_k}$ يو اټومي !!{formula} دے۔ له دې پرته
د !!{formula}s سټ هماغسې تعريفېږي لکه د لومړۍ درجې په منطق کښې، خو د
دويمې درجې کميت ټاکنې اضافي شرطونه لري۔ پام وکړئ چې !!{identity}
يوازې د لومړۍ درجې د ترمونو دپاره لرو: که $R$ او $S$ د يوه شان
ځايښت~$k$ اړيکيز !!{variable}s وي، $\eq[R][S]$ د لاندې عبارت د لنډون
په توګه تعريفولے شو:
\[
\lforall[x_1 \dots][\lforall[x_k][(\Atom{R}{x_1, \dots, x_k} \liff
  \Atom{S}{x_1, \dots, x_k})]].
\]

د دويمې درجې منطق قاعدې يوازې د کميت ټاکونکو قاعدې د دويمې درجې
نويو متغيرونو ته غځوي۔ خو دلته بايد لږ پام وشي چې دا متغيرونه د
ژبې~$\Lang{L}$ له !!{predicate}s، او په عام ډول د~$\Lang{L}$ له
!!{formula}s سره څنګه چلند کوي۔ لږ تر لږه، اړيکيز متغيرونه ترمونه
ګڼل کېږي، نو د لاندې بڼې استنتاجونه لرو:
\[
!A(R) \Proves \lexists[R][!A(R)]
\]
خو که $\Lang{L}$ د حساب ژبه وي او د ثابتې اړيکې نښه~$<$ ولري، تمه
کېږي چې لاندې استنتاج هم معتبر وي:
\[
x < y \Proves \lexists[R][\Atom{R}{x,y}]
\]
يا د يوه ورکړل شوي !!{formula}~$!A$ دپاره:
\[
!A(x_1, \dots, x_k) \Proves \lexists[R][\Atom{R}{x_1,\dots,x_k}]
\]
په لا عام ډول، ښايي د لاندې بڼې استنتاجونه روا کول وغواړو:
\[
\Subst{!A}{\lambd[\vec x][!B(\vec x)]}{R} \Proves
\lexists[R][!A]
\]
دلته $\Subst{!A}{\lambd[\vec x][!B(\vec x)]}{R}$ هغه پايله ښيي چې
په~$!A$ کښې د $\Atom{R}{t_1,\dots,t_k}$ بڼې هر اټومي !!{formula} په
$!B(t_1, \dots, t_k)$ بدل شي۔ دا وروستۍ قاعده د {\em ادراک شېما} له
شتون سره برابره ده، يعنې د لاندې بڼې يو بديهي اصل:
\[
\lexists[R][\lforall[x_1, \dots, x_k][(!A(x_1, \dots, x_k) \liff
\Atom{R}{x_1, \dots, x_k})]],
\]
د دويمې درجې په ژبه کښې د هر !!{formula}~$!A$ دپاره يوه داسې بېلګه
چې $R$ پکې ازاد متغير نۀ وي۔ (تمرین: وښيئ چې که $R$ ته اجازه ورکړل
شي چې په~$!A$ کښې راشي، دا شېما ناسازګاره ده!)
\textbf{د ژباړې د نسخې سمون (OLFOL-031):} د بدلولو په تشريح کښې
سرچينې د اړيکې اټومي نښه د څيز په نښه ليکلې وه؛ دلته د اټومي اړيکې
سمه نښه بېرته ايښوول شوې ده۔

کله چې منطقپوهان د ``دويمې درجې منطق بديهي اصول'' يادوي، په عام ډول
د دويمې درجې د کميت ټاکونکو پر قاعدو او د ادراک پر شېما د لومړۍ درجې
منطق لږ تر لږه غځونه مرادوي۔ خو د دې بديهي اصولو او قاعدو د کمزورو
فرعي نظامونو مطالعه هم ډېر ځله په زړه پورې وي۔ د بېلګې په توګه، پام
وکړئ چې د ادراک بديهي شېما په خپل بشپړ عموميت کښې
\emph{ناوړاندوينيزه} ده: دا د داسې اړيکې
$\Atom{R}{x_1, \dots, x_k}$ د شتون دعوه روا کوي چې په داسې
!!a{formula} ``تعريف'' شوې وي چې د دويمې درجې کميت ټاکونکي لري؛ او دا
کميت ټاکونکي د ټولو داسې اړيکو پر سټ ګرځي---هغه سټ چې خپله $R$ هم پکې
ده!{} د
شلمې پېړۍ د پيل شاوخوا د راسل تناقض ته يو عام غبرګون دا و چې دا ډول
تعريفونه ملامت وګڼل شي او د رياضياتو د بنسټونو په جوړولو کښې ترې ډډه
وشي۔ که په !!{formula}~$!A$ کښې د دويمې درجې کميت ټاکونکو کارونه منع
شي، د ادراک يوه {\em وړاندوينيزه} بڼه ترلاسه کېږي چې څه ناڅه کمزورې
ده۔

د معنايي نظر له مخې، د دويمې درجې يو !!{structure} داسې ګڼلے شو چې
د ژبې دپاره د لومړۍ درجې يو !!{structure} او ورسره پر !!{domain} د
هغو اړيکو يو سټ لري چې د دويمې درجې کميت ټاکونکي پرې ګرځي (په لا
دقيق ډول، د هر~$k$ دپاره د ځايښت~$k$ اړيکو يو سټ شته)۔ البته که
ادراک په !!{derivation} نظام کښې شامل وي، دا اضافي شرط لرو چې د
``دويمې درجې په برخه'' کښې د ادراک د بديهي اصولو د پوره کولو دپاره
کافي اړيکې وي---ګنې د !!{derivation} نظام سم نۀ دے!{} د کافي اړيکو د
تضمين يوه اسانه لار دا ده چې د دويمې درجې برخه د لومړۍ درجې پر برخه
له \emph{ټولو} اړيکو جوړه وګڼو۔ داسې !!a{structure} ته \emph{بشپړ}
ويل کېږي، او په يوه معنا په رښتيا د ژبې ``ټاکل شوے !!{structure}''
دے۔ که خپل پام بشپړو !!{structure}s ته محدود کړو، هغه څه ترلاسه کوو
چې د دويمې درجې \emph{بشپړه} معنٰی پوهنه بلل کېږي۔ په دې حالت کښې د
يوه !!{structure} مشخصول يوازې د لومړۍ درجې د برخې مشخصولو ته راکمېږي،
ځکه د دويمې درجې د برخې منځپانګه په ضمني ډول ترې معلومېږي۔

په لنډيز کښې، د دويمې درجې منطق په خبرو کښې څه ابهام شته۔ د
!!{derivation} نظام له مخې ښايي له لاندې دواړو څخه يو مراد وي:
\begin{enumerate}
\item د دويمې درجې يو ``لږ تر لږه'' !!{derivation} نظام، له ځينو
  ادراکي بديهي اصولو سره۔
\item د دويمې درجې ``معياري'' !!{derivation} نظام، له بشپړ ادراک سره۔
\end{enumerate}
د معنٰی پوهنې له مخې هم ښايي له لاندې دواړو څخه يوه د پام وړ وي:
\begin{enumerate}
\item ``کمزورې'' معنٰی پوهنه، چې !!a{structure} پکې د لومړۍ درجې يوه
  برخه او د ادراک د بديهي اصولو د پوره کولو دپاره کافي لويه د دويمې
  درجې برخه لري۔
\item د دويمې درجې ``معياري'' معنٰی پوهنه، چې يوازې بشپړ
  !!{structure}s پکې څېړل کېږي۔
\end{enumerate}
کله چې منطقپوهان خپل مراد !!{derivation} نظام يا معنٰی پوهنه نۀ
مشخصوي، اکثره د هر لړ دويم توکے مرادوي۔ لکه چې وبه وينو، د دې معنٰی
پوهنې ګټه دا ده چې د ډېرو طبيعي رياضيکي جوړښتونو قطعي توصيفونه راکوي؛
په ورته وخت کښې !!{derivation} نظام ډېر پياوړے او د دې معنٰی پوهنې
دپاره سم دے۔ نيمګړتيا ئې دا ده چې !!{derivation} نظام د دې معنٰی
پوهنې دپاره \emph{بشپړ نۀ} دے؛ په حقيقت کښې د دويمې درجې د بشپړې
معنٰی پوهنې دپاره \emph{هېڅ} په اغېزمن ډول ورکړل شوے !!{derivation}
نظام بشپړ نۀ دے۔ له بلې خوا به ووينو چې !!{derivation} نظام د کمزورې
معنٰی پوهنې دپاره \emph{بشپړ} دے؛ نو که يوه جمله د اثبات وړ نۀ وي،
يو \emph{داسې} !!{structure} شته، چې لازم نۀ دے بشپړ وي، او جمله پکې
دروغ ده۔

د دويمې درجې منطق ژبه ډېره بډايه ده۔ يوځاييزې اړيکې د !!{domain} له
فرعي سټونو سره يو شان ګڼلے شو، نو په ځانګړي ډول پر دغو سټونو کميت
ټاکلے شو؛ د بېلګې په توګه، د طبيعي عددونو دپاره استقرا په يوه بديهي
اصل بيانولے شو:
\[
\lforall[R][((\Atom{R}{\Obj{0}} \land \lforall[x][(\Atom{R}{x} \lif
    \Atom{R}{x'})]) \lif \lforall[x][\Atom{R}{x}])].
\]
که د حساب ژبه د $\Obj 0, \Obj \prime, +, \times$ او~$<$ نښې ولري، د
هغوى د چلند د بيانولو دپاره لاندې بديهي اصول ورزياتولے شو:
\begin{enumerate}
\item $\lforall[x][\lnot x' = \Obj 0]$
\item $\lforall[x][\lforall[y][(x' = y' \lif x = y)]]$
\item $\lforall[x][(x + \Obj 0) = x]$
\item $\lforall[x][\lforall[y][(x + y') = (x + y)']]$
\item $\lforall[x][(x \times \Obj 0) = \Obj 0]$
\item $\lforall[x][\lforall[y][(x \times y') = ((x \times y) + x)]]$
\item $\lforall[x][\lforall[y][(x < y \liff \lexists[z][y = (x + z')])]]$
\end{enumerate}
دا ښودل ګران نۀ دي چې دغه بديهي اصول د پاسني استقرا له بديهي اصل سره،
په دې شرط چې د دويمې درجې بشپړه معنٰی پوهنه کاروو، د حساب د معياري
!!{structure}~$\Struct{N}$، د حساب د معياري مدل، قطعي توصيف وړاندې
کوي۔ که کوم !!{structure}~$\Struct{M}$
دا بديهي اصول رښتوني کوي، له~$\Nat$ څخه د~$\Struct{M}$ !!{domain} ته
تابع~$f$ پر~$\Nat$ د عادي بازګښت په وسيله داسې تعريف کړئ چې
$f(0) = \Assign{\Obj 0}{M}$ او $f(x+1) = \Assign{\prime}{M}(f(x))$ وي۔
پر~$\Nat$ د عادي استقرا او په~$\Struct M$ کښې د بديهي اصولو (1)
او~(2) له صدق څخه وينو چې $f$ !!{injective} ده۔ د دې د ښودلو دپاره
چې $f$ !!{surjective} ده، $P$ د~$\Domain{M}$ د هغو غړو سټ ونيسئ چې
د~$f$ د قيمتونو په سټ کښې دي۔ ځکه چې $\Struct M$ بشپړ دے، $P$ د
دويمې درجې په !!{domain} کښې دے۔ د~$f$ له جوړونې پوهېږو چې
$\Assign{\Obj 0}{M}$ په~$P$ کښې دے او $P$ د~$\Assign{\prime}{M}$
لاندې تړلے دے۔ دا چې د استقرا بديهي اصل په~$\Struct M$ کښې (په ځانګړي
ډول د~$P$ دپاره) صدق کوي، تضمينوي چې $P$ د~$\Struct M$ د لومړۍ درجې
له ټول !!{domain} سره برابر دے۔ نو $f$ !!a{bijection} ده۔ دا
ښودل هم زيات ګران نۀ دي چې $f$ هومومورفيزم دے؛ پر~$\Nat$ عادي استقرا
بيا بيا کاروو۔
\textbf{د ژباړې د نسخې سمون (OLFOL-032):} په دويم بديهي اصل کښې
سرچينې د تابعې يوه بېله نښه کارولې وه، که څه هم ژبه او پاتې بحث د
پريـم نښه کاروي؛ دلته دويم اصل هماغې پريـم نښې ته يو شان شوے دے۔

د سټ تيورۍ په اصطلاحاتو کښې تابع يوازې د اړيکې يو ځانګړے ډول دے؛ د
بېلګې په توګه، يوځاييزه تابع~$f$ له داسې دوه ځاييزې اړيکې~$R$ سره يو
شان ګڼل کېدے شي چې $\lforall[x][\lexists![y][R(x,y)]]$ پوره کوي۔ نو
پر تابعو هم کميت ټاکلے شو۔ بيا د بشپړې معنٰی پوهنې په کارولو د
لامتناهي !!{structure}s ټولګے د هغو !!{structure}s~$\Struct M$ د
ټولګي په توګه تعريفولے شو چې د~$\Struct M$ له !!{domain} څخه د هغه
يو مناسب فرعي سټ ته يو پر يو تابع لري:
\[
\lexists[f][(\lforall[x][\lforall[y][(\eq[f(x)][f(y)] \lif \eq[x][y])]]
  \land \lexists[y][\lforall[x][\eq/[f(x)][y]]])].
\]
بيا د دې جملې نفي د متناهي !!{structure}s ټولګے تعريفوي۔

سربېره پر دې، د خطي ترتيب تعريف ته د لاندې شرط په زياتولو د ښو
ترتيبونو ټولګے تعريفولے شو:
\[
\lforall[P][(\lexists[x][\Atom{P}{x}] \lif \lexists[x][(\Atom{P}{x}
    \land \lforall[y][(y < x \lif \lnot \Atom{P}{y})])])].
\]
دا وايي چې هر ناتش سټ تر ټولو کوچنے غړے لري، که ``سټ'' له ``يوځاييزې
اړيکې'' سره يو شان وګڼو۔ د بلې بېلګې دپاره، د ګرافونو تړلتيا داسې
بيانولے شو چې د څوکو هېڅ ناتشه بېلتون په ناتړلو برخو کښې نشته:
\[
\lnot \lexists[A][(\lexists[x][A(x)] \land \lexists[y][\lnot A(y)]
  \land \lforall[w][\lforall[z][((\Atom{A}{w} \land \lnot \Atom{A}{z})
      \lif \lnot \Atom{R}{w,z})]])].
\]
د يوې بلې بېلګې په توګه، په تمرين کښې هڅه کولے شئ د هغو متناهي
!!{structure}s ټولګے تعريف کړئ چې !!{domain} ئې جفته اندازه لري۔ تر
دې هم په زړه پورې دا چې حقيقي عددونه د ناطقو عددونو لرونکي بشپړ مرتب
ميدان په توګه په قطعي ډول توصيفولے شو۔

په لنډ ډول، د دويمې درجې منطق د لومړۍ درجې له منطق څخه ډېر زيات
بيانوونکے دے۔ دا ښه خبر و؛ اوس بد خبر ته راځو۔ مخکې مو وويل چې د
دويمې درجې د بشپړې معنٰی پوهنې دپاره هېڅ بشپړ اغېزمن !!{derivation}
نظام نشته۔ که ښه ئې وګڼو که بد، د لومړۍ درجې منطق ډېر خاصيتونه دلته
نشته، چې د متناهي شاهد خاصيت او د لوېنهايم--سکولم قضيې هم پکې دي۔

له بلې خوا، که څوک د دويمې درجې بشپړه معنٰی پوهنه پرېږدي او کمزورې
بڼه ئې ومني، نو د دويمې درجې لږ تر لږه !!{derivation} نظام د دې معنٰی
پوهنې دپاره بشپړ دے۔ په نورو ټکو، که $\Proves$ داسې ولولو چې ``په لږ
تر لږه نظام کښې ثابتوي'' او $\Entails$ داسې چې ``په کمزورې معنٰی پوهنه
کښې معنايي استلزام لري''، ښودلے شو چې هر کله $\Gamma \Entails !A$ وي
نو $\Gamma \Proves !A$ وي۔ که څوک په !!{derivation} نظام کښې ځانګړي
ادراکي بديهي اصول شاملول غواړي، بايد معنٰی پوهنه د دويمې درجې هغو
جوړښتونو ته محدوده کړي چې دا اصول پوره کوي: د بېلګې په توګه، که
$\Delta$ د ادراکي بديهي اصولو يو سټ وي (کېدے شي ټول وي)، نو که
$\Gamma \cup \Delta \Entails !A$ وي، $\Gamma \cup \Delta \Proves !A$
هم وي۔ په ځانګړي ډول، که $!A$ د هغو ادراکي بديهي اصولو په کارولو د
اثبات وړ نۀ وي چې موږ ئې څېړو، د~$\lnot !A$ داسې يو مدل شته چې بيا هم
دغه ادراکي بديهي اصول پوره کوي۔

دا چې د بشپړتيا قضيه د کمزورې معنٰی پوهنې دپاره ولې صدق کوي، تر ټولو
اسانه ئې دا ده چې د دويمې درجې منطق په لاندې ډول د ډېرو ډولونو منطق
وګڼو۔ يو ډول د عادي ``لومړۍ درجې'' !!{domain} په توګه تفسيرېږي، او
بيا د هر~$k$ دپاره د ``ځايښت~$k$ اړيکو'' !!a{domain} لرو۔ ژبې ته
جوړې دننه اړيکيزې نښې ``$\Atom{\Obj{true}_k}{R,x_1,\dots,x_k}$''
ورکوو؛ مراد ئې دا دے چې $R$ د~$x_1$، \dots،~$x_k$ دپاره صدق کوي،
چرته چې $R$ د ``$k$-ځاييزې اړيکې'' د ډول متغير او $x_1$،
\dots،~$x_k$ د لومړۍ درجې د ډول څيزونه دي۔

د دې يو شان ګڼلو له مخې، د دويمې درجې کمزورې معنٰی پوهنه په اصل کښې د
ډېرو ډولونو د منطق عادي معنٰی پوهنه ده؛ او مخکې مو ليدلي چې د ډېرو
ډولونو منطق د لومړۍ درجې په منطق کښې ځايېدلے شي۔ نو د دواړو خواوو
ژباړو ته په پام، د دويمې درجې منطق کمزورے تصور په رښتيا د لومړۍ درجې
منطق يوه پټه بڼه ده چې !!{domain} ئې هم ``څيزونه'' او هم ``اړيکې''
لري او مناسب بديهي اصول پرې واکمن دي۔

\end{document}
